% ===================================================================== DISCUSSION
\section{Discussion}
\label{sec:discussion}

\paragraph{Engine scope}
The proof applies to an engine pair only when both satisfy the interface and adapter
obligations of \cref{def:engine,tab:adapter}: deterministic finality and state
roots, side-effect-free shadow validation, root-seeded target initialization,
and durable source retirement. The PoA$\to$HotStuff path is the evaluated
partial-path instantiation. A crash-tolerant Raft~\cite{raft} source adapter
with majority replication and internal fence rejection crosses a finalized
boundary in a component test; this is not a second distributed migration
campaign. Its source quorum changes the admission condition: at $n=3f+1$,
$q_1=\lfloor n/2\rfloor+1$ satisfies $q_1>2f$ only when $n\ge4f$.
Thus a majority Raft source at $(7,2)$ is not admitted without a stronger
source policy. The target-safety premise of \cref{th:safety} is independent
of HotStuff internals, but conforming Bullshark~\cite{bullshark},
Mysticeti~\cite{mysticeti}, or HotStuff-1~\cite{hotstuff1} adapters and their
execution evidence have not been supplied.

\paragraph{Retirement timing}
Under durable Rule~\ref{rule:lock}, a valid $n-f$ \CutCert{} and admitted
$q_1>2f$ already exclude a residual source quorum. The subsequent \FenceCert{}
records retirement bound to the manifest and admitted policy; it is not a
second independent necessity for source quiescence in that locked design.
The result in \cref{prop:final-parent-retirement} offers a conditional alternative
to that retirement timing: universal final-parent mediation and atomic stop-on-Boundary
adoption could retire the source. It does not replace Rule~\ref{rule:lock}. The proposition
does not apply to speculative or pipelined source engines that may authorize a
child before its parent finalizes. Its source serialization cost and progress
are not measured by the existing campaigns; neither later target activation
nor any current campaign interpretation follows from this alternative.

\paragraph{Composition with autonomous switching policies}
A learned or threat-adaptive policy~\cite{adachain,threatadapt,adacon}
could \emph{propose} a migration schedule, but governance would still have to
admit the source policy, ordered heights, fixed committee, and every premise
of the boundary proof. The handoff guarantee is conditional on the admitted
schedule; it does not validate the policy that selected it. Within one
governance epoch, the \CutCert{} and terminal locks constrain authority.
No cooldown limits how often governance can open another epoch, so migration
rate-limiting remains a deployment responsibility.

\paragraph{Client semantics}
A client that acts on a provisional block may later see its suffix discarded.
GRANDPA's optimistic-versus-final distinction and the confirmation-depth
convention of deployed chains offer analogous choices, but neither supplies
receipt-revocation semantics~\cite{grandpa}. The signed tier says whether a receipt is revocable,
not what to undo if it is revoked; receipt reconciliation and client compensation remain outside
the live path. Irreversible cross-chain effects must wait for a
\emph{verified} \SealCert{} covering that receipt; observing $h\ge\hr$ alone
does not make provisional effects final.

\paragraph{Evidence required for a system claim}
The primary missing experiment is a networked forward handoff: distribute
\CutCert{} and \ManCert{} shares, durably adopt the manifest and install local
fences, form and distribute \FenceCert, activate the target, and complete
\emph{both} certified abort and seal branches. Committee-wide fence deployment is not
implemented in the evaluated network path; manifest gossip and terminal-certificate dissemination
are absent from that path. Component predicates and the event-driven controller specification
are not an execution of this transport (\cref{tab:status}). Crash/restart and partial-delivery
injection must bracket
each durable write and certificate transition, including a first post-restart
source action. A timeout or a gate that safely refuses transfer is not a
successful migration. The existing restart campaigns show catch-up in particular
crash-only schedules, not persistent enforcement of Rule~\ref{rule:lock}
throughout the handoff. Bounded TLC checks cover selected sizes, and equality
of finite model/code outcome sets is not a prefix-preserving trace refinement;
the faithful $n=7$ partial-delivery run has no completed verdict.

The inadmissible, quorum-confounded RQ1 cost data in \cref{sec:rq1} cannot
become full-path latency evidence by normalization or tighter confidence
intervals. A rerun must hold the relevant quorum policies fixed across
comparable arms, use an admitted source, collect all certificates, and record
per-trial service interruption and client-completed throughput. To establish
comparative systems performance, a pinned artifact-backed switch (for example,
BFTBrain~\cite{bftbrain} or a Bedrock-based construction~\cite{bedrock})
and a Fabric-style maintenance halt must be executed on shared hardware,
network, workload, and client-completion definitions. Only common dimensions
can be compared: the cited comparators were not executed, an internal Cox-style
arm is not Cox, and no competitive external-client benchmark is available.

The serialized simulator cannot realize a concurrent two-leader fork because
it uses one proposer per height. Loopback partition controls are deliberately
constructed, and the three-trial single-host impairment grid cannot estimate
operational risk. Retained same-region recovery hosts add narrow catch-up evidence,
not physical geo-WAN behaviour: physical geo-WAN remains outside this evaluation.
Configured-endpoint throughput labels do not independently authenticate placement.
Larger campaigns need per-trial authority traces and campaign-specific host,
impairment, toolchain, and timing records. The current rollback fixtures
check context admission, not execution-level replay and root equivalence
under real application/oracle dependencies. Beyond the empirical gaps, the
proof assumes one fixed identical committee and monolithic state; joint
membership change, sharding, safe close-and-retry of stale candidates, and
client compensation for discarded provisional effects remain outside it.

\paragraph{Cost and feasibility}
The retained synthetic evaluation needs no external dataset. The canonical
single-host Ed25519 aggregate has a median of about $3.9$\,ms for $n=200$
(134 shares; \cref{tab:certtiming}); individual timing samples are not retained.
Neither panel measures distributed quorum latency: the share-list CPU result omits
transport, while the simulator cutover-evidence panel omits committee-wide dissemination.

Canonical-record accounting gives about $119$--$158$\,KiB
of replay context per 1000 provisional blocks, excluding application replay
and external reconciliation. Historical whole-command intervals are available
for three input families, but missing campaign hardware and per-trial timing
prevent a defensible wall-clock or monetary estimate for the entire evaluation.

The successful design intentionally pauses source voting for certification
before target activation. Its conditional bound in \eqref{eq:tbdy} includes
readiness, certificate collection/delivery and failover, durable writes,
initialization, and the first target commit; it does not guarantee a bound
when these premises fail. The retained $10$--$16$\,ms partial-path gaps omit
these stages. Neither the complete pause nor externally observed service
interruption has yet been measured.
